% Options for packages loaded elsewhere
\PassOptionsToPackage{unicode}{hyperref}
\PassOptionsToPackage{hyphens}{url}
\documentclass[
  ignorenonframetext,
]{beamer}
\newif\ifbibliography
\usepackage{pgfpages}
\setbeamertemplate{caption}[numbered]
\setbeamertemplate{caption label separator}{: }
\setbeamercolor{caption name}{fg=normal text.fg}
\beamertemplatenavigationsymbolsempty
% remove section numbering
\setbeamertemplate{part page}{
  \centering
  \begin{beamercolorbox}[sep=16pt,center]{part title}
    \usebeamerfont{part title}\insertpart\par
  \end{beamercolorbox}
}
\setbeamertemplate{section page}{
  \centering
  \begin{beamercolorbox}[sep=12pt,center]{section title}
    \usebeamerfont{section title}\insertsection\par
  \end{beamercolorbox}
}
\setbeamertemplate{subsection page}{
  \centering
  \begin{beamercolorbox}[sep=8pt,center]{subsection title}
    \usebeamerfont{subsection title}\insertsubsection\par
  \end{beamercolorbox}
}
% Prevent slide breaks in the middle of a paragraph
\widowpenalties 1 10000
\raggedbottom
\AtBeginPart{
  \frame{\partpage}
}
\AtBeginSection{
  \ifbibliography
  \else
    \frame{\sectionpage}
  \fi
}
\AtBeginSubsection{
  \frame{\subsectionpage}
}
\usepackage{iftex}
\ifPDFTeX
  \usepackage[T1]{fontenc}
  \usepackage[utf8]{inputenc}
  \usepackage{textcomp} % provide euro and other symbols
\else % if luatex or xetex
  \usepackage{unicode-math} % this also loads fontspec
  \defaultfontfeatures{Scale=MatchLowercase}
  \defaultfontfeatures[\rmfamily]{Ligatures=TeX,Scale=1}
\fi
\usepackage{lmodern}
\ifPDFTeX\else
  % xetex/luatex font selection
\fi
% Use upquote if available, for straight quotes in verbatim environments
\IfFileExists{upquote.sty}{\usepackage{upquote}}{}
\IfFileExists{microtype.sty}{% use microtype if available
  \usepackage[]{microtype}
  \UseMicrotypeSet[protrusion]{basicmath} % disable protrusion for tt fonts
}{}
\makeatletter
\@ifundefined{KOMAClassName}{% if non-KOMA class
  \IfFileExists{parskip.sty}{%
    \usepackage{parskip}
  }{% else
    \setlength{\parindent}{0pt}
    \setlength{\parskip}{6pt plus 2pt minus 1pt}}
}{% if KOMA class
  \KOMAoptions{parskip=half}}
\makeatother
\setlength{\emergencystretch}{3em} % prevent overfull lines
\providecommand{\tightlist}{%
  \setlength{\itemsep}{0pt}\setlength{\parskip}{0pt}}
\usepackage{bookmark}
\IfFileExists{xurl.sty}{\usepackage{xurl}}{} % add URL line breaks if available
\urlstyle{same}
\hypersetup{
  pdftitle={Abstract},
  hidelinks,
  pdfcreator={LaTeX via pandoc}}

\title{\texorpdfstring{Abstract}{Abstract}}
\author{\texorpdfstring{}{}}
\date{}

\begin{document}
\frame{\titlepage}

\begin{frame}{Abstract}
\protect\phantomsection\label{abstract}
This review unifies four converging research programmes---case systems
in theoretical linguistics, categorial grammar and its
string-diagrammatic semantics, enriched and topos-theoretic category
theory, and active inference as a process theory of cognition---into a
single coherent framework centred on the \emph{cognitive case diagram}.
The formal core treats case systems as categories (objects = case roles,
morphisms = grammatical relations), cross-linguistic alignment types as
functors between case categories, and sentences as compact-closed string
diagrams in the DisCoCat (Distributional Compositional Categorical)
framework. By demonstrating that DisCoCat constitutes the algebraic
formalization of the distributional programme that modern large language
models implement empirically---from static word embeddings through
contextual transformers---we bridge the two traditions of meaning
(formal and distributional) within a unified categorical architecture.
These diagrams serve simultaneously as mathematical proofs, cognitive
representations, and the structural core of generative models under
active inference: understanding a sentence reduces to constructing the
case diagram that minimizes variational free energy. We extend this
framework through Distributional Active Inference, which models full
distributions over case diagrams rather than point estimates, connecting
distributional semantics to distributional reinforcement learning via
enriched-categorical push-forward measures. Graded proto-role structure
is captured via \([0,1]\)-enriched categories with categorical magnitude
as a complexity measure, and inter-theoretic translation is formalized
through Morita-equivalent classifying toposes. Implementation within the
CEREBRUM (Case-Enabled Reasoning Engine with Bayesian Representations
for Unified Models) architecture demonstrates practical viability. The
analysis suggests that the mathematics of case alignment is not merely
an abstraction of human grammar, but the formal geometry of relational
meaning itself---the architecture through which cognitive systems
organize experience into who-does-what-to-whom.
\end{frame}

\end{document}
